\documentclass[12pt]{book}
 \usepackage[a4paper,headheight=16pt,scale={0.8,0.85},hoffset=0.4cm]{geometry}

% Para escribir en castellano:
\usepackage[spanish]{babel}
\usepackage[utf8]{inputenc}

\usepackage[pdftex]{graphicx}
\usepackage{amsmath}
\usepackage{amssymb}
\usepackage{tipa}
%\usepackage{txfonts}
\usepackage{textcomp}
%\usepackage{amsthm}
%\usepackage{array}
%\usepackage{xy}
\usepackage{fancyhdr}

\pagestyle{fancy}
\renewcommand{\chaptermark}[1]{\markboth{#1}{}}
\renewcommand{\sectionmark}[1]{\markright{\thesection\ #1}}
%Eliminando numeración y nombres de Capítulos y secciones
\addto\captionsspanish{
  \renewcommand\chaptername{}
}
\renewcommand{\thechapter}{}
\setcounter{secnumdepth}{-1} 


\fancyhf{}
\fancyhead[LE,RO]{\bfseries\thepage}
\fancyhead[LO]{\bfseries\rightmark}
\fancyhead[RE]{\bfseries\leftmark}
\fancyfoot[R]{\includegraphics[height=1cm]{footer_logo}}
\renewcommand{\headrulewidth}{0.5pt}
\renewcommand{\footrulewidth}{0pt}
\addtolength{\headheight}{0.5pt}
\setlength{\footskip}{0.5in}
\renewcommand{\footruleskip}{0pt}
\fancypagestyle{plain}{%
 \fancyhead{}
 \renewcommand{\headrulewidth}{0pt}
}

%
%\parindent 0in
\parskip 0.05in

\usepackage{hyperref}
\hypersetup{
            pdftitle=Proyecto Nahual – Introducción a la programación,%
            pdfauthor=Proyecto Nahual,%
            pdfsubject=Libro de Introducción a la programación del Proyecto Nahual,%
            pdfkeywords=programación nahual ruby,%
            colorlinks,%
            citecolor=black,%
            filecolor=black,%
            linkcolor=black,%
            urlcolor=black,%
            pdftex}

%%%%%%%%%%%%%%%%%%%%%%%%%%%%%%%%%%%%%%%%%%%%%%%%%%%%%%%%%%%%%%%%%%%%%%%%%%%%%%%%
\usepackage{color}
\definecolor{gray75}{gray}{.75}

\definecolor{listinggray}{gray}{0.9}
\definecolor{lbcolor}{rgb}{0.9,0.9,0.9}
\usepackage{listings}
\usepackage{textcomp}
\lstset{
	backgroundcolor=\color{lbcolor},
	tabsize=4,
	rulecolor=,
	language=ruby,
        basicstyle=\scriptsize,
        upquote=true,
        aboveskip={1.5\baselineskip},
        columns=fixed,
        showstringspaces=false,
        extendedchars=true,
        breaklines=true,
        prebreak = \raisebox{0ex}[0ex][0ex]{\ensuremath{\hookleftarrow}},
        frame=single,
        showtabs=false,
        showspaces=false,
        showstringspaces=false,
        identifierstyle=\ttfamily,
        keywordstyle=\color[rgb]{0,0,1},
        commentstyle=\color[rgb]{0.133,0.545,0.133},
        stringstyle=\color[rgb]{0.627,0.126,0.941},
}
%Ejemplo de uso para incluir un archivo: \lstinputlisting{parameters.py}
%Ejemplo para codigo inline: \begin{lstlisting} \end{lstlisting}

\newlength{\PictHOffset}
\newlength{\PictVOffset}
\setlength{\PictHOffset}{1in}
\addtolength{\PictHOffset}{\hoffset}
\addtolength{\PictHOffset}{\oddsidemargin}

\setlength{\PictVOffset}{1in}
\addtolength{\PictVOffset}{\voffset}
\addtolength{\PictVOffset}{\topmargin}
\addtolength{\PictVOffset}{\headheight}
\addtolength{\PictVOffset}{\headsep}
\addtolength{\PictVOffset}{\topskip}
\addtolength{\PictVOffset}{-1\paperheight}
\begin{document}
\begin{titlepage}
   \noindent\hspace*{-\PictHOffset}%
   % Uncomment next line and replace XXX by your graphics file
   \raisebox{\PictVOffset}[0pt][0pt]{\makebox[0pt][l]{%
      \includegraphics[width=\paperwidth,height=\paperheight]{portada}}}
   % Comment next line.
    %\raisebox{\PictVOffset}[0pt][0pt]{\makebox[0pt][l]{\color{blue}\rule{\paperwidth}{0.25\paperheight}}}

   \begin{center}
      \vfill
      \bfseries \Huge My Master Thesis \\[3cm]
      by Myself
      \vfill
     University of \LaTeX land
   \end{center}
\end{titlepage}

\thispagestyle{empty}
\subsection*{\center \normalsize Alguna cosa sobre el libro}
 
\tableofcontents 
\mainmatter
\chapter{1era semana - ¿Qué es programar?}
\section{¿Qué es programar?}
Recuerdo claramente el momento en que escribí mi primer programa (y eso que yo tengo una memoria malísima), era un programa bastante complicado para alguien como yo. Pero como yo estudiaba matemática en la facultad se suponía que era bueno en cosas como ``pensar lógicamente''. Entonces fui a la sala de computadoras de la facu, armado solamente con un libro sobre programación y mi ego, me senté en una de las computadoras y empecé. Bueno, ``empecé'' no sería la palabra correcta. O tampoco ``programar". Digamos que me senté delante de la computadora y me sentí increíblemente tonto. Luego me sentí avergonzado. Y luego enojado hasta que finalmente me sentí pequeño: después de 8 horas desesperantes el programa estaba terminado y andando. El hecho de que funcionase no importaba mucho...era un momento de gloria.\\

Pero la pregunta que me quedo fue, ¿por qué fue tan difícil decirle a una computadora que haga algo medianamente complicado? Bueno, la parte difícil no es la de ``medianamente complicado'' sino la parte de ``decirle a una computadora``.\\

En cualquier comunicación entre humanos uno puede dejar muchas cosas sin decir, saltearse pasos y conceptos y entenderse. Claro que uno deja que la otra persona ``complete'' lo que falta. En realidad, uno tiene que saltearse cosas. Si no lo hiciésemos, ¡no podríamos hacer nada! El ejemplo típico es el  de armarse un sandwich de jamón, queso y mayonesa. Normalmente, uno diría tranquilamente: ``Che ¿me haces un sandwich de jamón y queso con mayonesa?'' En cambio, si se lo tenés que pedir a alguien que nunca hizo uno, como un robot, tendrías que decirle algo como:
\begin{enumerate}
  \item Agarrá 2 rodajas de pan (y dejá el resto)
  \item Sacá queso, jamón, mayonesa y un cuchillo
  \item Untá un poco de mayonesa en un pan
  \item Poné arriba de la mayonesa: 2 fetas de jamón y 1 de queso
  \item Poné la otra rodaja de pan arriba, poné todo en un plato y traémelo. ¡Gracias!
\end{enumerate}

Me imagino que esas instrucciones son suficientes para un niño de 7 años. Pero un niño de 7 años es bastante inteligente como para hacer eso. En cambio ¿qué habría que decirle a una computadora? Concentrémonos solamente en el primer paso:

\begin{enumerate}
  \item Agarrá 2 rodajas de pan (y dejá el resto)
  \begin{enumerate}
    \item Localizá la bolsa de pan lactal
    \item Andá a una mesa
    \item Dejá la bolsa de pan en la mesa
    \item Agarrá 2 rodajas de adentro de la bolsa
    \item \ldots
  \end{enumerate}
\end{enumerate}

Lamentablemente estas instrucciones no alcanzan para nada. Por ejemplo: ¿cómo es “localizá la bolsa de pan”? Tendríamos que tener una especie de base de datos en donde se asocien las cosas con su ubicación. La base de datos debería también tener las coordenadas del jamón, heladera, mesa, cuchillo, plato, yo mismo...
¡Ah! ¿Y qué pasaría si la bolsa de pan está en una panera? Entonces ¿habría que abrir la panera? O ¿por ahí está detrás de algo? ¿O arriba de la heladera? Entonces tendríamos solamente para localizar la bolsa de pan lactal:
\begin{enumerate}
  \item Buscá las coordenadas de la panera
  \item Si la bolsa de pan está en la panera
  \begin{enumerate}
    \item Abrí la panera
    \item Meté la mano dentro de la panera
    \item Agarrá la bolsa de pan lactal
    \item Sacá la mano de la panera
    \item Cerrá la panera
  \end{enumerate}
  \item Si la bolsa de pan está arriba de la heladera
  \begin{enumerate}
    \item Andá hasta la puerta de la heladera
    \item Estirá la mano hacia arriba de la heladera
    \item Agarrá la bolsa de pan lactal
    \item Sacá la mano de arriba de la heladera
  \end{enumerate}
\end{enumerate}

Y así sigue y sigue. Además ¿qué pasa si no hay cuchillo? ¿O si la mesa está llena de cosas y no hay lugar para preparar el sandwich? ¡Habría que rezar al Santo que te parezca más apropiado para que el pan esté en buen estado! Aún los pasos más sencillos, como abrir la panera, deberían ser explicados ¡y es por eso que todavía no hay robots que hagan sandwich para nosotros! No es que los robots no puedan ser construidos ¡sino que no podemos programarlos para que nos hagan un sandwich! Y eso es porque, aunque no lo parezca, preparar un sandwich de jamón, queso y mayonesa es algo muy difícil de describir (pero fácil de hacer para criaturas inteligentes como nosotros, los humanos) y las computadoras sólo son buenas para las cosas que son fáciles de describir. Y por eso para mí fue tan difícil hacer mi primer programa: ¡las computadoras son más tontas de lo que creía!

\section{Hacer un sandwich es programar}
Cuando uno le explica a alguien como hacer un sandwich, el trabajo es mucho más fácil porque esa persona ya sabe lo que es un sandwich. Es común que la gente sepa de “sandwichería“ y por eso esa persona puede llenar los espacios de nuestra pobre explicación. El paso 3 dice “Untá un poco de mayonesa en un pan“, pero no dice que no lo unte en los dos lados. Tampoco dice cuánta mayonesa ni de qué manera untarla (si poner más en el medio, o en los costados) ni aclara que debe usar el cuchillo para hacerlo y no las manos. Uno asume que esas cosas ya se saben. Creo que esta realidad ayudaría a explicar qué quiere decir programar o por lo menos dar una definición informal.\\

Programar es decirle a tu computadora cómo hacer algo. Las tareas grandes tienen que ser divididas o descompuestas en tareas más chicas, las cuales tienen que ser descompuestas en tareas aún más chicas y así seguir hasta que las tareas sean tan chiquititas y simples que no tienen que ser divididas porque la computadora sabe hacerlas.\\

Estas tareas son cosas muy básicas como imprimir un texto en la pantalla o sumar 2 números. El problema más grande que tuve cuando estaba aprendiendo a programar fue que yo estaba aprendiendo al revés. Yo ya sabía lo que quería que la computadora hiciese y trataba de ir dividiendo eso hasta llegar a tareas que de la computadora entendiese. Muy mala idea. En realidad yo no sabía lo que realmente la computadora podía hacer, entonces no sabía bien en que tareítas descomponer mi tarea grande o lo que la computadora tenía que hacer.\\
 	
Por eso en Nahual, vas a aprender a programar en forma diferente. Primero vas a aprender estas tareítas básicas que la computadora puede hacer (algunas nomás) y después vas a encontrar tareas simples que puedan ser descompuestas en esas tareítas. Tu primer programa va a ser tan fácil que solamente te llevara menos de un minuto.

\section{Lenguajes de Programación}
Para poder decirle algo a la computadora tenemos que usar un lenguaje diferente. Un lenguaje de programación es parecido al lenguaje que usamos para hablar, en el sentido que está formado por cosas básicas como palabras (verbos y sustantivos) y por maneras de combinar esas palabras para dar sentido (oraciones, párrafos, novelas). Hay muchos lenguajes para elegir (C, Java, Ruby, Perl, \ldots) y cada uno tiene su conjunto de elementos básicos. Algunos tienen más que otros. Ruby tiene un conjunto sencillo pero elegante y es uno de los más fáciles para aprender así que vamos a usar ese.\\
 
Probablemente la mejor razón para usar Ruby es que los programas en Ruby suelen ser cortos. Por ejemplo aquí hay un pequeño programa en Java.

\begin{lstlisting}[language=java]
public class Hola Mundo {
  public static void main (String [] args) {
    System.out.println("Hola Mundo");
  }
}
\end{lstlisting}

Y acá esta el mismo programa pero en el lenguaje Ruby:

\begin{lstlisting}
  puts 'Hola Mundo'
\end{lstlisting}

Este programa, en ambos casos, escribe las palabras ``Hola Mundo'' en la pantalla de la computadora. En la versión en Ruby uno podría adivinar, pero en la versión de java adivinar esto no es tan sencillo.

\section{El arte de programar}
Una parte importante de programar es, claramente, que el programa haga lo que se supone que tiene que hacer y no otra cosa. En otras palabras, el programa no debería tener errores o \emph{bugs} (se pronuncia bags). Claro que, cuando uno se concentra mucho en no tener errores, se pierde la gran mayoría de las cosas buenas de la programación.\\

Programar no es solamente acerca del producto final o del programa sino acerca del proceso que lleva a él. Los programas, a diferencia de cosas más concretas como una casa, no se construyen en un orden tan rígido sino que los programadores charlan sobre él, piensan, discuten, escriben un poco, lo arreglan, lo estiran, juegan con él, lo cambian, lo mejoran, lo prueban, lo borran, lo vuelven a escribir, \ldots

\section{Un programa no se arma: crece}
Dado que un programa siempre está creciendo y evolucionando, siempre está cambiando y tiene que estar escrito pensando en que va a cambiar. Sé que esto no puede quedar muy claro hablando en términos prácticos pero este tema ira surgiendo durante todo el libro. Probablemente, la primera regla de un buen programador es evitar duplicar código, esto es ``No vuelva a escribir lo mismo". Si la vemos de otra manera, es como la ley del mínimo esfuerzo. Un buen programador tiene que ser un poco vago, pero no ser vago de no hacer nada, sino ser vago en forma inteligente ¡para no tener que hacer nada después! ¡Ahórrate trabajo siempre que puedas! Si hacer cambios ahora significa ahorrarte trabajo más adelante ¡hacelo!\\

Trata de que tu programa haga lo menos posible para hacer lo que tiene que hacer. De esta manera programar no solo es mucho más interesante (porque es muy aburrido hacer lo mismo una y otra vez) sino que, además ¡el programa tiene menos errores y es más rápido! ¡Es una situación donde todos ganan!\\

De cualquier manera, la idea es la misma: tratá de que tus programas sean flexibles. Así, cuando haya que hacer los cambios (y te aseguro que siempre vienen) entonces vas a pasar menos y mejor tiempo haciéndolos.

Ejercicio 1 - Mostrar un programa
Les mostramos un programa hecho en Ruby de los robots.
En este capítulo proponemos un role-play con un robot. Los chicos tienen que escribir los pasos del robot y el robot (uno de los instructores) tiene que repetirlo.
\chapter{2da semana - Números y letras}
\nota{Introducción}{En la primer semana aprendimos qué era programar, vimos un robot en funcionamiento y tratamos de ver los conceptos de instrucciones en ese ejemplo}
\nota{Objetivo}{El paso que queremos dar en esta semana es empezar a hacer esto pero con la computadora y con el lenguaje de programación Ruby}
\nota{Cómo lo vamos a hacer}{Vamos a empezar por el tema de los números que es algo bastante sencillo para arrancar y que se usa durante toda la vida del programador, de hecho es muy probable que vuelvan sobre esto en el futuro}

Cuando terminemos con los números vamos a pasar a cómo mostrar las palabras por la pantalla y, por supuesto en el medio vamos a tratar de jugar, crear y pensar; que, de hecho, es lo que más nos interesa que hagamos juntos

\section{Números}
Ahora que ya tenemos todo listo, abran su editor de texto y escriban:

\begin{lstlisting}
  puts 1 + 2
\end{lstlisting}

Guarden el programa (con CTRL + S) y pónganle como nombre \emph{calculadora.rb}. Ahora ejecuten el programa con \emph{ruby calculadora.rb} en la línea de comando (en la ventana de DOS) y debería poner un 3 en la pantalla. ¿Vieron? ¡No fue tan difícil!

\section{Introducción a \emph{puts}}
Entonces ¿qué está pasando en ese programa? Estoy seguro de que podemos adivinar qué es lo que está pasando con el 1+2: nuestro programa hace lo mismo que este:

\begin{lstlisting}
  puts 3
\end{lstlisting}

La instrucción puts sencillamente imprime en la pantalla aquello que le pones después, por eso el puts 3 va a dejar en el monitor que tienen enfrente un número 3.

\section{Enteros y con coma}
En la mayoría de los lenguajes de programación (y Ruby no es la excepción) los números sin coma se llaman enteros y los números con coma se llaman de punto flotante o ¡con coma! Por ejemplo, acá les mostramos algunos números enteros: 

\begin{itemize}
  \item 5
  \item -205
  \item 99999999999999999999
  \item 0
\end{itemize}

Y acá van algunos con punto flotante: 

\begin{itemize}
  \item 54.321
  \item 0.001
  \item -205.3884
  \item 0.0
\end{itemize}

En la práctica no solemos usar  números con coma, solamente usamos números enteros (son los más comunes). Claro que hay una razón para ello y es que las cosas de todos los días se nos presentan de manera entera: tenemos 2 hermanos, alquilamos 4 películas, tenemos 17 años, etc. Los números con coma los usamos cuando queremos ser mucho más precisos y en general se usan para cosas científicas locas. Hasta la mayoría de los programas que manejan plata no usan la coma ¡llevan las monedas como enteros!

\section{Vamos a hacer cuentas}
Hasta ahora, tenemos todo como para hacer nuestra calculadora. Las calculadoras usan números con coma así que si queremos que nuestro programa haga de calculadora ¡vamos a tener que usarlos! Los números los escribimos con el teclado (están en la parte de arriba). Para la coma, en Ruby se usa el . (Punto) y no la coma.  Esto se debe a que los que hicieron Ruby viven en Japón, donde no se usa la coma para los decimales sino el Punto(.)\\

Para la suma y la resta usamos los clásicos signos + y - , como vimos en el primer ejemplo. Para la multiplicación usamos el asterisco (*) y para la división usamos la barra (/). Algunos teclados tienen todos estos símbolos a la derecha.\\

Ahora que sabemos esto, vamos a modificar el programa \emph{calculadora.rb} para extenderlo un poco más. 

\begin{lstlisting}
puts 1.0 + 2.0
puts 2.0 * 3.0
puts 5.0 - 8.0
puts 9.0 / 2.0
\end{lstlisting}

Esto es lo que nuestro programa escribiría por pantalla (acordate que \emph{puts} sirve para escribir algo en la pantalla, en este caso escribe el resultado de una cuenta):
\begin{console-output}
3.0
6.0
-3.0
4.5
\end{console-output}

Ahora vamos a probar con números enteros, o sea sin los decimales:

\begin{lstlisting}
puts 1+2
puts 2*3
puts 5-8
puts 9/2
\end{lstlisting}

Los resultados (lo que el programa escribe en la pantalla) son casi los mismos ¿no?...

\begin{console-output}
3
6
-3
4
\end{console-output}

... Salvo por ese 4 al final. Lo que sucede es que cuando la computadora hace cuentas con números enteros, el resultado siempre va a ser un entero. Y si  al hacer una división, el resultado no es un número entero sino que tiene decimales, entonces redondea para abajo. Pero ¿para qué sirve la división entera?\\

Imaginemos que la entrada al cine sale \$8 y tenemos \$50. ¿Cuántas entradas podemos comprar con \$50? ¡6! No podemos comprar 6,25. Digo, con los \$2 que sobran de vuelto ¡en el cine no te van a dar un cuarto de entrada!. Algunas cosas no son divisibles.\\

Ahora vamos a experimentar un poco. Si queremos hacer cuentas más difíciles, podemos usar los paréntesis... por ejemplo:

\begin{lstlisting}
puts 5 * (12-8) + -15
puts 98 + (59872 / (13*8)) * -51
\end{lstlisting}

Lo que da como resultado
\begin{console-output}
5
-29227
\end{console-output}

\section{Letras y Palabras}
Ahora que ya aprendimos sobre los números en Ruby, vamos a ver qué pasa con las letras ¿y las palabras? En el mundo de la programación solemos referirnos a las letras como caracteres y a las palabras como cadenas de caracteres o strings (en inglés). En Ruby, las cadenas de caracteres siempre van encerradas entre ' ' (comillas simples) o “ ” (comillas dobles). Ejemplos de cadenas de caracteres serían:

\begin{itemize}
  \item 'Hola.'
  \item 'Ruby esta re-bueno'
  \item '5 es mi número favorito...y el tuyo?'
  \item "Cuando me martillo un dedo digo \^?*!@."
  \item ' '
\end{itemize}

Como podemos ver, algunas cadenas no son solo con letras del abecedario. Los caracteres también pueden ser los signos de puntuación, los números, los espacios, el signo de pregunta y cualquier simbolito que puedas tener en el teclado. Si la cadena no tiene caracteres entonces se la denomina cadena vacía.\\

Si recuerdan, habíamos usado la instrucción \emph{puts} para escribir por la pantalla el resultado de una cuenta o un número. ¡Ahora la vamos a usar para escribir cadenas! Para esto hagamos un nuevo programa que se llamará \emph{presentacion.rb}

\nota{Aclaración}{Todos los programas de Ruby tienen la extensión .rb por ese motivo ponemos esto al final de cada uno de los nombres de los programas}

\begin{lstlisting}
  puts 'Hola a todos'
  puts 'mi nombre es Charly'
  puts 'uf.........()\%/?'
\end{lstlisting}

¡Ustedes fíjense que es lo que el programa hace ahora!

\section{Operaciones con las cadenas}

Con las cadenas también podemos hacer algunas cuentas locas. No todas, pero algunas si...por ejemplo:

\begin{lstlisting}
  puts 'me gusta' + 'el apio'
\end{lstlisting}

Para poder ejecutar este programa hay que escribir \emph{ruby presentacion.rb} y en la ventana de DOS vemos:
\begin{console-output}
  me gustael apio
\end{console-output}

¡Uy! ¡Pucha! ¡Me olvide de poner un espacio! ¡No voy a esperar que Ruby sepa castellano también! Hay un dicho que dice: “Las computadoras no van a hacer lo que vos querés, van a hacer lo que les dijiste que hagan”. Para corregir esto podemos hacer cualquiera de estas tres opciones:

\begin{lstlisting}
  puts 'me gusta ' + 'el apio'
  puts 'me gusta' + ' ' + 'el apio'
  puts 'me gusta' + ' el apio'
\end{lstlisting}

Como se puede ver (si ejecutamos el programa) no es muy importante donde ponemos el espacio. Fíjense que después del + hay un espacio en el segundo ejemplo y ninguno en el primero.\\%FIXME: No entiendo esta nota

Ahora sabemos que podemos sumar cadenas, pero ¿qué pasa si hacemos?:

\begin{lstlisting}
  puts 'hola soy charly'*4
\end{lstlisting}

El programa debería responder:
\begin{console-output}
  ya te escuche... cuantas veces me lo vas a decir
\end{console-output}

¡Ja! ¡No! ¡Ruby no es taaaaaan piola! Lo que se debería ver es:

\begin{console-output}
hola soy charlyhola soy charlyhola soy charlyhola soy charly
\end{console-output}

Bueno, vamos a jugar un poco con todo esto ahora:
\begin{lstlisting}
  puts 12 + 12
  puts ' 12' + ' 12'
  puts ' 12 + 12'
\end{lstlisting}

Esto daría:
\begin{console-output}
  24
  1212
  12 + 12
\end{console-output}

En el primer caso el resultado da 24 porque Ruby entiende que suma dos números enteros, sin embargo en el segundo caso se da cuenta de que son dos cadenas y por eso las concatena, pero en el último caso, lo único que le estamos diciendo que muestre por pantalla '12 + 12' como si fuera una palabra. ¿Y qué tal esto?:
\begin{lstlisting}
  puts 2 * 5
  puts ' 2' * 5
  puts ' 2 * 5'
\end{lstlisting}

A ver...
\begin{console-output}
10
22222
2 * 5
\end{console-output}

\section{Ya empezamos a tener problemas...}
A esta altura, capaz que alguno se sintió tentado de probar alguna cosa que no funciona, como:
\begin{lstlisting}
  puts '12' + 12
  puts '2' * '5'
\end{lstlisting}

Esto da un error como:
\begin{console-output}
  #<TypeError: cannot convert Fixnum into String>
\end{console-output}

Lo que dice Ruby es algo así como “No puedo mezclar peras con bananas”. En realidad, lo que dice es que no puede convertir un número en una cadena de caracteres. El verdadero problema es que no se puede sumar el 12 al '12' ya que sería sumar un número a un texto. La línea de abajo tampoco funcionaría porque no podemos multiplicar dos cadenas. Simplemente ¡Ruby no te deja! (además, hacer eso sería ilógico)
 
\begin{thebibliography}{99}
  
\bibitem{algo}
Titulo de la pagina (autor o algo): \\
http://url.com
 
\bibitem{otro}
Autor, \emph{Titulo}, Editorial (año) \\
ISBN: XX-XXXX-XXX-X
 
\end{thebibliography}
 
\end{document}